\begin{theorem}[Absorbing edge gauges into the second PEPS]%
    \label{thm:peps_edge_gauge_absorption}
    \lean{TNLean.PEPS.edge_gauge_absorption}
    \leanok
    \uses{def:peps_absorbEdgeGauges, thm:peps_absorbEdgeGauges_bondDim,
        thm:peps_absorbEdgeGauges_component}
    Given edge gauges $X_e$, set $\widetilde B=B^X$. Then
    $D_{\widetilde B}=D_B$ and, for every vertex $v$,
    \begin{align}
        \widetilde B_v(\eta,\sigma)
        &=
        \sum_{\eta'}
        \left(\prod_{ie}M_{ie}(\eta(ie),\eta'(ie))\right)
        B_v(\eta',\sigma),
        \notag
    \end{align}
    where $M_{ie}=X_e$ or $M_{ie}=(X_e^{-1})^t$ according to the orientation of
    the endpoint $ie$.
    % Formula: the preceding component identity is gauge absorption at v.
    % Ink-to-index/type verdict: paper_normal.tex:1212--1314 gives only
    % vertex-specific pictures, each with boundary (V,P)=(deg(v),1); the old
    % degree-four star cannot represent the theorem on an arbitrary graph.
    After this absorption, arbitrary insertions of the same matrix on the same
    edge are compared in the first PEPS and the modified second PEPS.
\end{theorem}
\begin{proof}\leanok
    With $\widetilde B=B^X$, Theorem~\ref{thm:peps_absorbEdgeGauges_bondDim}
    gives $D_{\widetilde B}=D_B$, and
    Theorem~\ref{thm:peps_absorbEdgeGauges_component} gives
    \begin{align}
        \widetilde B_v(\eta,\sigma)
        &=
        \sum_{\eta'}
        \left(\prod_{ie}M_{ie}(\eta(ie),\eta'(ie))\right)
        B_v(\eta',\sigma).
        \notag
    \end{align}
\end{proof}

\begin{definition}[Post-absorption inserted-edge comparison]%
    \label{def:peps_postAbsorptionEdgeInsertionEquality}
    \lean{TNLean.PEPS.PostAbsorptionEdgeInsertionEquality}
    \leanok
    \uses{def:peps_edgeInsertedCoeff}
    For tensor families $A,\widetilde B$, the post-absorption comparison holds
    when $h:D_A=D_{\widetilde B}$ and the following equality holds. For each
    edge $e$, write
    $\phi_{h_e}:\MN{D_A(e)}
        \to\MN{D_{\widetilde B}(e)}$ for the matrix-algebra
    transport induced by $h_e:D_A(e)=D_{\widetilde B}(e)$. Then, for every
    matrix $X\in\MN{D_A(e)}$ and every physical configuration $\sigma$,
    $C_A(e,X;\sigma)
        =C_{\widetilde B}(e,\phi_{h_e}(X);\sigma)$.
\end{definition}

\begin{theorem}[Post-absorption edge insertion equality]%
    \label{thm:peps_postAbsorptionEdgeInsertionEquality}
    \lean{TNLean.PEPS.post_absorption_edge_insertion_equality}
    \leanok
    \uses{thm:peps_edge_gauge_absorption,
        def:peps_postAbsorptionEdgeInsertionEquality}
    Let $A$ and $B$ be vertex-injective PEPS with the same state, assume every
    bond dimension is positive, and assume the bond-dimension equality
    $h:D_A=D_B$. After the edge gauges have been
    absorbed into the second tensor family, there are gauges $Z_e$ and
    $\widetilde B=B^Z$ such that $D_{\widetilde B}=D_A$ and, for every edge
    $e$, write $\phi_{h_e}:\MN{D_A(e)}\to\MN{D_{\widetilde B}(e)}$ for the
    matrix-algebra transport induced by the corresponding equality
    $h_e:D_A(e)=D_{\widetilde B}(e)$. Then, for every matrix
    $X\in\MN{D_A(e)}$ and every physical configuration $\sigma$,
    $C_A(e,X;\sigma)
        =C_{\widetilde B}(e,\phi_{h_e}(X);\sigma)$.
    On the five-site graph used in the proof, this equality reads
    % Formula: C_A(e,X;sigma)=C_{\widetilde B}(e,phi_{h_e}(X);sigma), with
    % the transport identified on the common bond in the displayed house spelling.
    % Ink-to-index: each panel has the five physical tensors of eq:inj_equal_edge,
    % six graph bonds (the 5-cycle and the 2--5 chord), and X splits that chord.
    % All virtual indices are contracted; the five north legs are sigma_1,...,sigma_5.
    % Boundary signature: both source panels have (V,P)=(0,5), specializing the
    % general theorem's (0,|V(G)|) physical boundary.  Source verdict: exact graph
    % and insertion from paper_normal.tex:1037--1065.
    \begin{center}
        \tenkzkernel
        \begin{tenkz}[rows={wire}, cols=5, frame=circle, bonds=none]
            \tn[at=(1,1), name=edgeEqA1, label pos=145,
                ports={0:virtual, 180:virtual, 90:physical}]{A_1}
            \tn[at=(1,2), name=edgeEqA2, label pos=145,
                ports={0:virtual, 180:virtual, -72:virtual, 90:physical}]{A_2}
            \tn[at=(1,3), name=edgeEqA3, label pos=145,
                ports={0:virtual, 180:virtual, 90:physical}]{A_3}
            \tn[at=(1,4), name=edgeEqA4, label pos=145,
                ports={0:virtual, 180:virtual, 90:physical}]{A_4}
            \tn[at=(1,5), name=edgeEqA5, label pos=145,
                ports={0:virtual, 180:virtual, 252:virtual, 90:physical}]{A_5}
            \tnwire{edgeEqA1.0}{edgeEqA2.180}
            \tnwire{edgeEqA2.0}{edgeEqA3.180}
            \tnwire{edgeEqA3.0}{edgeEqA4.180}
            \tnwire{edgeEqA4.0}{edgeEqA5.180}
            \tnwire{edgeEqA5.0}{edgeEqA1.180}
            \tnwire[name=edgeEqAChord]{edgeEqA2.-72}{edgeEqA5.252}
            \tn[skin=ring, species=operator, at=on edgeEqAChord 0.5]{X}
        \end{tenkz}
        $ = $
        \begin{tenkz}[rows={wire}, cols=5, frame=circle, bonds=none]
            \tn[at=(1,1), name=edgeEqB1, label pos=145,
                ports={0:virtual, 180:virtual, 90:physical}]{\widetilde B_1}
            \tn[at=(1,2), name=edgeEqB2, label pos=145,
                ports={0:virtual, 180:virtual, -72:virtual, 90:physical}]{\widetilde B_2}
            \tn[at=(1,3), name=edgeEqB3, label pos=145,
                ports={0:virtual, 180:virtual, 90:physical}]{\widetilde B_3}
            \tn[at=(1,4), name=edgeEqB4, label pos=145,
                ports={0:virtual, 180:virtual, 90:physical}]{\widetilde B_4}
            \tn[at=(1,5), name=edgeEqB5, label pos=145,
                ports={0:virtual, 180:virtual, 252:virtual, 90:physical}]{\widetilde B_5}
            \tnwire{edgeEqB1.0}{edgeEqB2.180}
            \tnwire{edgeEqB2.0}{edgeEqB3.180}
            \tnwire{edgeEqB3.0}{edgeEqB4.180}
            \tnwire{edgeEqB4.0}{edgeEqB5.180}
            \tnwire{edgeEqB5.0}{edgeEqB1.180}
            \tnwire[name=edgeEqBChord]{edgeEqB2.-72}{edgeEqB5.252}
            \tn[skin=ring, species=operator, at=on edgeEqBChord 0.5]{X}
        \end{tenkz}
    \end{center}
    as in \cite[Section~3, lines 1037--1065]{Molnar2018NormalPEPS}. The
    positive-bond hypothesis is the source's standing assumption that injective
    PEPS have nonzero virtual bond spaces; it is needed because the edge gauges
    come from blocking the PEPS around each edge into a three-site injective
    chain.
\end{theorem}
\begin{proof}\leanok
    \uses{thm:peps_exists_edgeGaugeFamily, thm:peps_edge_gauge_absorption,
        def:peps_postAbsorptionEdgeInsertionEquality}
    Block the PEPS around each edge into a three-site injective chain and compare
    the matched matrix insertions: this assigns to every edge an invertible bond
    matrix $X_e$ realizing the insertion-algebra transport as conjugation. Absorb
    the transpose of $X_e$ into the second tensor family on that edge. The
    open-edge gauge action then conjugates each inserted matrix by the transpose
    of the absorbed gauge, while the insertion-algebra transport conjugates by
    $X_e$; the two conjugations agree because the absorbed gauge is the transpose
    of $X_e$, giving the post-absorption insertion equality.
\end{proof}

\begin{theorem}[Local gauge extraction]%
    \label{thm:localGauge_exists}
    \lean{TNLean.PEPS.localGauge_exists}
    \lean{TNLean.PEPS.localGauge_exists_of_factorizedLocalGauge}
    \leanok
    \uses{def:peps_hasFactorizedLocalGauge}
    Under the factorized local gauge hypothesis, the local gauge formula follows.
    What remains open is to upgrade equality of the edge-blocked coefficients to
    the blocked-middle formula using the corresponding three-site
    injective-chain argument. The physical-realization and physical-to-virtual
    steps in that argument are recorded separately as
    Theorems~\ref{thm:peps_virtualInsertionPhysicalRealization} and
    \ref{thm:peps_physical_to_virtual_insertion}.
    After Theorem~\ref{thm:peps_postAbsorptionEdgeInsertionEquality} supplies
    the factorized local gauge hypothesis, injectivity on the star neighborhood
    provides a left inverse that isolates the local edge gauges. Thus there are
    incident-edge matrices $X_e$ such that
    % Formula: local extraction yields the valence-independent component equation.
    % Ink-to-index/type verdict: both sides have boundary (V,P)=(deg(v),1).
    % No standalone source picture states this extraction, and the old two panels
    % had unequal virtual boundary types, so the exact equation replaces them.
    \begin{align}
        B_v(\eta,\sigma)
        &=
        \sum_{\eta'}
        \left(\prod_{e\ni v}X_e(\eta_e,\eta'_e)\right)
        A_v(\eta',\sigma).
        \notag
    \end{align}
\end{theorem}
\begin{proof}\leanok
    The factorized local gauge datum already supplies the incident-edge
    invertible matrices and the factorized local coefficient identity.
\end{proof}

\begin{theorem}[Gauge consistency]%
    \label{thm:gaugeConsistency}
    \lean{TNLean.PEPS.gaugeConsistency}
    \leanok
    \uses{thm:peps_exists_edgeGaugeFamily,
        thm:peps_postAbsorptionEdgeInsertionEquality,
        thm:localGauge_exists, def:gaugeVertex}
    Suppose $A$ and $B$ are vertex-injective PEPS tensors with the same PEPS
    state, that the corresponding virtual bond dimensions are equal, that
    every virtual bond of $A$ has positive dimension, and that the graph is
    connected. The connectivity assumption is necessary: on a disconnected
    graph the per-vertex scalars produced by the reduction below cannot be
    absorbed into edge gauges, since on the empty graph on two vertices the
    products of the vertex scalars can agree while no edge gauge relates the two
    tensors.
    Choose an edge $e=(u,v)$ and block all other vertices into one middle
    tensor. The three resulting tensors form an injective chain, so the
    one-dimensional injective comparison assigns an invertible matrix to the
    edge $e$. Repeating this for every edge and absorbing these matrices into
    the second tensor family gives the post-absorption insertion equality of
    \cite[Section~3]{Molnar2018NormalPEPS}: for every edge and every matrix,
    inserting that matrix on the edge in the first PEPS gives the same state as
    inserting it on the same edge in the modified second PEPS. Applying the
    local extraction at each vertex then gives one edge gauge family whose
    endpoint actions transform the first tensor family into the second:
    % Formula: one global family X supplies the local gauge equation at every v.
    % Ink-to-index/type verdict: the local component has boundary
    % (V,P)=(deg(v),1), while a two-vertex comparison has
    % (deg(u)+deg(w)-2,2).  The source gives the vertexwise equations in
    % paper_normal.tex:1212--1314, not the leg-suppressed scalar legacy picture.
    \begin{align}
        B_v(\eta,\sigma)
        &=
        (A^X)_v(\eta,\sigma)
        =
        \sum_{\eta'}
        \left(\prod_{e\ni v}X_{v,e}(\eta_e,\eta'_e)\right)
        A_v(\eta',\sigma).
        \notag
    \end{align}
\end{theorem}
\begin{proof}\leanok
    Absorbing the edge gauges into the second tensor family gives the
    post-absorption insertion equality. On each vertex with an incident edge,
    the one-vertex-versus-complement comparison supplies a nonzero scalar
    relating the first tensor to the absorbed second tensor. Because the state
    is nonzero, these per-vertex scalars multiply to one, so on the connected
    graph a spanning-tree construction realizes their reciprocals as the
    oriented incidence product of one edge-scalar family. Folding those scalars
    into the inverse of the absorbed edge gauges produces the global gauge. On a
    one-vertex graph there are no edges and the gauge is trivial.
\end{proof}

\begin{remark}
    The preceding identities are the algebraic content of the edge-gauge
    construction in \cite[Section~3]{Molnar2018NormalPEPS}. On an ordered edge
    $e=(u,w)$, the actions $X_e$ at $u$ and $(X_e^{-1})^\top$ at $w$ cancel
    upon contraction. Absorbing these matrices into $B$ gives
    $\widetilde B$, for which every inserted-edge coefficient agrees with the
    corresponding coefficient of $A$. Comparing one vertex with its injective
    complement then yields the local gauge equation. Repeating the argument on
    every edge produces a single compatible gauge family; the two-injective
    comparison below shows that the remaining three-leg operator is scalar.
\end{remark}

\begin{definition}[Left-identity three-leg residual form]%
    \label{def:peps_threeLegLeftIdentityForm}
    \lean{TNLean.PEPS.threeLegLeftIdentityForm}
    \leanok
    For three virtual leg spaces $\alpha,\beta,\gamma$ and an endomorphism
    $Z$ of $\beta\otimes\gamma$, define the induced endomorphism of
    $\alpha\otimes\beta\otimes\gamma$ by acting as the identity on the
    $\alpha$ leg and as $Z$ on the remaining two legs. This is the
    $\Id_\alpha\otimes Z$ form appearing after the first tensor is inverted in
    \cite[Section~3]{Molnar2018NormalPEPS}.
\end{definition}

\begin{definition}[Right-identity three-leg residual form]%
    \label{def:peps_threeLegRightIdentityForm}
    \lean{TNLean.PEPS.threeLegRightIdentityForm}
    \leanok
    For an endomorphism $U$ of $\alpha\otimes\beta$, define the induced
    endomorphism of $\alpha\otimes\beta\otimes\gamma$ by acting as $U$ on the
    first two legs and as the identity on the $\gamma$ leg. This is the
    $U\otimes\Id_\gamma$ form in \cite[Section~3]{Molnar2018NormalPEPS}.
\end{definition}

\begin{definition}[Middle-identity three-leg residual form]%
    \label{def:peps_threeLegMiddleIdentityForm}
    \lean{TNLean.PEPS.threeLegMiddleIdentityForm}
    \leanok
    For an endomorphism $W$ of $\alpha\otimes\gamma$, define the induced
    endomorphism of $\alpha\otimes\beta\otimes\gamma$ by acting as $W$ on the
    outer two legs and as the identity on the middle $\beta$ leg. This is the
    third residual form in \cite[Section~3]{Molnar2018NormalPEPS}.
\end{definition}

\begin{theorem}[Scalar reduction of the residual two-tensor gauges]%
    \label{thm:peps_twoInjectiveGaugeScalarReduction}
    \lean{TNLean.PEPS.threeLeg_residual_forms_scalar}
    \leanok
    \uses{def:peps_threeLegLeftIdentityForm,
        def:peps_threeLegRightIdentityForm,
        def:peps_threeLegMiddleIdentityForm}
    In the proof of the generalized two-injective-tensor comparison, applying
    the inverse of the second injective tensor leaves three possible virtual
    operators on the exposed legs of the first tensor:
    % Formula (paper_normal.tex:1157--1193): A_1 equals B_1 with a two-leg
    % residual Z, equals B_1 with U on the complementary pair, and equals B_1
    % with W on the first and third legs.  Ink-to-index: each operator box has
    % two incoming and two outgoing virtual ports; the untouched third index is
    % the remaining direct wire.  The north leg is the common physical index.
    % Boundary signature: every panel has (V,P)=(3,1).  Source verdict: exact
    % four-panel scalar-reduction identity in house style.
    \begin{center}
        \tenkzkernel
        \begin{tenkz}[rows={wire}, cols=1, bonds=none]
            \tn[name=scalarA, label pos=w,
                ports={35:virtual, 0:virtual, -35:virtual,
                       90:physical}]{A_1}
        \end{tenkz}
        $ = $
        \begin{tenkz}[rows={wire}, cols=2, bonds=none]
            \tn[name=scalarBZ, label pos=w,
                ports={35:virtual, 0:virtual, -35:virtual,
                       90:physical}]{B_1}
            \tn[at=0.4 s of 1.2 e of scalarBZ, name=scalarZ, skin=box,
                species=operator,
                ports={160:virtual, 200:virtual, 20:virtual,
                       -20:virtual}]{Z}
            \tnwire{scalarBZ.0}{scalarZ.160}
            \tnwire{scalarBZ.-35}{scalarZ.200}
        \end{tenkz}
        $ = $
        \begin{tenkz}[rows={wire}, cols=2, bonds=none]
            \tn[name=scalarBU, label pos=w,
                ports={35:virtual, 0:virtual, -35:virtual,
                       90:physical}]{B_1}
            \tn[at=0.4 n of 1.2 e of scalarBU, name=scalarU, skin=box,
                species=operator,
                ports={160:virtual, 200:virtual, 20:virtual,
                       -20:virtual}]{U}
            \tnwire{scalarBU.35}{scalarU.160}
            \tnwire{scalarBU.0}{scalarU.200}
        \end{tenkz}
        $ = $
        \begin{tenkz}[rows={wire}, cols=2, bonds=none]
            \tn[name=scalarBW, label pos=w,
                ports={35:virtual, 0:virtual, -35:virtual,
                       90:physical}]{B_1}
            \tn[at=0.7 n of 1.2 e of scalarBW, name=scalarW, skin=box,
                species=operator,
                ports={160:virtual, 200:virtual, 20:virtual,
                       -20:virtual}]{W}
            \tnwire{scalarBW.35}{scalarW.160}
            \tnwire{scalarBW.-35}{scalarW.200}
        \end{tenkz}
    \end{center}
    Inverting the first injective tensor then shows that the corresponding
    residual operator on the three exposed virtual legs has the three
    decompositions
    \begin{align}
        \sum_i \Id\otimes Z_i^{(1)}\otimes Z_i^{(2)}
        &=
        \sum_i U_i^{(1)}\otimes U_i^{(2)}\otimes\Id
        =
        \sum_i W_i^{(1)}\otimes\Id\otimes W_i^{(2)}.
        \notag
    \end{align}
    Equivalently, for nonzero virtual leg spaces $\alpha,\beta,\gamma$, if one
    endomorphism of $\alpha\otimes\beta\otimes\gamma$ lies simultaneously in
    the three subspaces $\Id_\alpha\otimes\End(\beta\otimes\gamma)$,
    $\End(\alpha\otimes\beta)\otimes\Id_\gamma$, and the analogous
    subspace with $\Id_\beta$ in the middle leg, then it is a scalar multiple of
    the identity. Thus the residual virtual operators are scalar multiples of
    the identity. This is the scalar step in the proof of
    \cite[Section~3]{Molnar2018NormalPEPS}.
\end{theorem}
\begin{proof}\leanok
    Write the common operator as $T$. The equality
    $T=\Id_\alpha\otimes Z=U\otimes\Id_\gamma$ gives
    $\delta_{aa'}Z_{bc,b'c'}=U_{ab,a'b'}\delta_{cc'}$ for all indices.
    If $c\neq c'$, the right hand side is zero; choosing $a=a'$ gives
    $Z_{bc,b'c'}=0$. Hence $Z$ is diagonal in the $\gamma$-index. Similarly,
    the equality $T=\Id_\alpha\otimes Z=W_{\alpha\gamma}\otimes\Id_\beta$
    gives
    $\delta_{aa'}Z_{bc,b'c'}=W_{ac,a'c'}\delta_{bb'}$,
    so $Z_{bc,b'c'}=0$ whenever $b\neq b'$. Thus $Z_{bc,b'c'}=0$ unless
    $b=b'$ and $c=c'$.

    It remains to compare the surviving diagonal entries. From
    $\delta_{aa'}Z_{bc,bc}=U_{ab,a'b}$, the value
    $Z_{bc,bc}$ is independent of $c$. From the middle-identity equality it is
    also independent of $b$. Therefore
    $Z=\lambda\Id_{\beta\otimes\gamma}$. Substitution into
    \begin{align}
        \Id_\alpha\otimes Z
        &=
        U\otimes\Id_\gamma
        =
        W_{\alpha\gamma}\otimes\Id_\beta
        \notag
    \end{align}
    gives $U=\lambda\Id_{\alpha\otimes\beta}$ and
    $W_{\alpha\gamma}=\lambda\Id_{\alpha\otimes\gamma}$.
\end{proof}

\begin{theorem}[Many-leg scalar extraction]%
    \label{thm:peps_piProductFormsScalar}
    \lean{TNLean.PEPS.piProduct_forms_scalar}
    \leanok
    Let a finite family of legs be given, each a nonzero finite index set, with
    invertible matrices $M_i$ and $N_i$ on each leg. If
    $\prod_i(M_i)_{p_iq_i}=\prod_i(N_i)_{p_iq_i}$ for every configuration pair
    $(p,q)$, then there are nonzero scalars $c_i$ with $M_i=c_iN_i$ for each
    leg and $\prod_i c_i=1$.
\end{theorem}
\begin{proof}\leanok
    Each $N_i$ is invertible, so each of its rows has a nonzero entry; pick a
    reference configuration $(p^0,q^0)$ on which every $(N_i)_{p^0_iq^0_i}$ is
    nonzero. The product identity at $(p^0,q^0)$ then makes every
    $(M_i)_{p^0_iq^0_i}$ nonzero as well. Fix a leg $i$ and a pair $(a,b)$, and
    apply the identity to the configuration that agrees with the reference
    except at leg $i$, where it takes $(a,b)$. Splitting off leg $i$ gives
    $(M_i)_{ab}K_i=(N_i)_{ab}L_i$, where $K_i,L_i$ are the products of the
    other legs at the reference, both nonzero. The reference identity at leg
    $i$ relates $K_i$ and $L_i$, and eliminating them yields
    $(M_i)_{ab}=c_i(N_i)_{ab}$ with
    $c_i=(M_i)_{p^0_iq^0_i}/(N_i)_{p^0_iq^0_i}$, independent of $(a,b)$.
    Taking the product of the scalars over all legs and using the reference
    identity gives $\prod_i c_i=1$.
\end{proof}

\begin{theorem}[Pointwise product cancellation for two injective blocks]%
    \label{thm:peps_twoBlockPointwiseProductCancellation}
    \lean{TNLean.PEPS.twoBlockReciprocalScalarProportional_of_pointwise_mul_eq}
    \leanok
    Let $A_1,A_2,B_1,B_2$ be tensors with matching external,
    shared-boundary, and physical index sets. Assume that the two external
    boundary spaces and the shared-boundary configuration space are nonempty,
    and that $A_1$ and $A_2$ are injective. Suppose that, for every pair of
    external indices, every pair of shared-boundary configurations, and every
    pair of physical indices,
    \begin{align}
        A_1(\eta_1,\mu,\sigma_1)A_2(\eta_2,\nu,\sigma_2)
        &=
        B_1(\eta_1,\mu,\sigma_1)B_2(\eta_2,\nu,\sigma_2).
        \notag
    \end{align}
    Then there exists $\lambda\in\C^\times$ such that
    $A_1=\lambda B_1$ and $A_2=\lambda^{-1}B_2$.
\end{theorem}
% This theorem is the rank-one cancellation step once
% \cite[Section~3]{Molnar2018NormalPEPS} has been reduced to pointwise product
% equality of separated coefficients.
\begin{proof}\leanok
    Choose external and shared-boundary indices, and use injectivity of
    $A_1$ and $A_2$ to choose physical indices for which the two chosen
    coefficients are nonzero. The product equality then shows that the
    corresponding coefficients of $B_1$ and $B_2$ are also nonzero. Define
    $\lambda=
        B_2(\eta_2^0,\nu^0,\sigma_2^0)/
        A_2(\eta_2^0,\nu^0,\sigma_2^0)$.
    For every $(\eta_1,\mu,\sigma_1)$, the product identity with
    $(\eta_2^0,\nu^0,\sigma_2^0)$ fixed gives
    \begin{align}
        A_1(\eta_1,\mu,\sigma_1)
        A_2(\eta_2^0,\nu^0,\sigma_2^0)
        &=
        B_1(\eta_1,\mu,\sigma_1)
        B_2(\eta_2^0,\nu^0,\sigma_2^0).
        \notag
    \end{align}
    Dividing by $A_2(\eta_2^0,\nu^0,\sigma_2^0)\neq0$ gives
    $A_1=\lambda B_1$. Fixing $(\eta_1^0,\mu^0,\sigma_1^0)$ and substituting
    $A_1=\lambda B_1$ into the product identity gives
    \begin{align}
        \lambda B_1(\eta_1^0,\mu^0,\sigma_1^0)
        A_2(\eta_2,\nu,\sigma_2)
        &=
        B_1(\eta_1^0,\mu^0,\sigma_1^0)
        B_2(\eta_2,\nu,\sigma_2).
        \notag
    \end{align}
    Dividing by $B_1(\eta_1^0,\mu^0,\sigma_1^0)\neq0$ gives
    $A_2=\lambda^{-1}B_2$.
\end{proof}

\begin{definition}[Matrix unit for coefficient extraction]%
    \label{def:peps_matrixUnit}
    \lean{TNLean.PEPS.matrixUnit}
    \leanok
    For indices $i$ and $j$, let $E_{ij}$ denote the matrix whose only
    nonzero entry is the $(i,j)$ entry, where it is equal to $1$.
\end{definition}

\begin{theorem}[Coefficient extraction for one shared bond]%
    \label{thm:peps_twoBlockInsertedCoeff_singletonBond}
    \lean{TNLean.PEPS.twoBlockInsertedCoeff_singletonBond_single}
    \leanok
    \uses{def:peps_matrixUnit}
    Suppose that the two blocks share a single virtual bond. Then inserting
    the matrix unit $E_{\mu,\nu}$ in that bond extracts the corresponding
    product of tensor coefficients:
    \begin{align}
        C_{A_1,A_2}
        (b,E_{\mu,\nu};\eta_1,\eta_2,\sigma_1,\sigma_2)
        &=
        A_1(\eta_1,\mu,\sigma_1)A_2(\eta_2,\nu,\sigma_2).
        \notag
    \end{align}
\end{theorem}
\begin{proof}\leanok
    Since there is only one shared bond, the condition that the two shared
    configurations agree away from the distinguished bond is vacuous, and
    \begin{align}
        C_{A_1,A_2}
        (b,E_{\mu,\nu};\eta_1,\eta_2,\sigma_1,\sigma_2)
        &=
        \sum_{\mu',\nu'}
        (E_{\mu,\nu})_{\mu'(b),\nu'(b)}
        A_1(\eta_1,\mu',\sigma_1)A_2(\eta_2,\nu',\sigma_2).
        \notag
    \end{align}
    The entries of the matrix unit satisfy
    $(E_{\mu,\nu})_{i,j}
        =\delta_{i,\mu(b)}\delta_{j,\nu(b)}$.
    Thus the only nonzero summand is the one with
    $(\mu',\nu')=(\mu,\nu)$.
\end{proof}

\begin{theorem}[One-bond two-injective comparison]%
    \label{thm:peps_twoInjectiveTensorInsertionComparison_singletonBond}
    \lean{TNLean.PEPS.two_injective_tensor_insertion_comparison_singletonBond}
    \leanok
    \uses{thm:peps_twoBlockInsertedCoeff_singletonBond,
        thm:peps_twoBlockPointwiseProductCancellation}
    Suppose that the two injective blocks have exactly one shared virtual
    bond and nonempty external boundary spaces. If all matrix insertions on
    that bond give equal two-block coefficients for the pairs
    $(A_1,A_2)$ and $(B_1,B_2)$, then there exists
    $\lambda\in\C^\times$ such that
    $A_1=\lambda B_1$ and $A_2=\lambda^{-1}B_2$.
\end{theorem}
\begin{proof}\leanok
    Insert the matrix unit $E_{\mu,\nu}$ and use the coefficient-extraction
    theorem to obtain the pointwise product equality
    \begin{align}
        A_1(\eta_1,\mu,\sigma_1)A_2(\eta_2,\nu,\sigma_2)
        &=
        B_1(\eta_1,\mu,\sigma_1)B_2(\eta_2,\nu,\sigma_2)
        \notag
    \end{align}
    for all indices. The pointwise product cancellation theorem gives the
    reciprocal scalar proportionality.
\end{proof}

\begin{theorem}[Generalized two-injective-tensor comparison]%
    \label{thm:peps_twoInjectiveTensorInsertionComparison}
    \lean{TNLean.PEPS.two_injective_tensor_insertion_comparison}
    \lean{TNLean.PEPS.SameAwayFromBond, TNLean.PEPS.twoBlockInsertedCoeff}
    \lean{TNLean.PEPS.IsTwoBlockInjective, TNLean.PEPS.SameTwoBlockInsertions}
    \lean{TNLean.PEPS.TwoBlockScalarProportional}
    \lean{TNLean.PEPS.TwoBlockReciprocalScalarProportional}
    \leanok
    \uses{thm:peps_twoInjectiveGaugeScalarReduction}
    Let $A_1,A_2,B_1,B_2$ be injective tensors with the same shared virtual
    bonds and nonempty external boundary spaces, and assume that there is at
    least one shared virtual bond. For the set $\mathcal B$ of shared bonds, a
    bond $b\in\mathcal B$, and $X\in\MN{D_b}$, set
    \begin{align}
        C_{A_1,A_2}(b,X;\eta_1,\eta_2,\sigma_1,\sigma_2)
        &=
        \sum_{\substack{\mu,\nu\\
        \mu|_{\mathcal B\setminus\{b\}}
        =\nu|_{\mathcal B\setminus\{b\}}}}
        X_{\mu_b,\nu_b}
        A_1(\eta_1,\mu,\sigma_1)A_2(\eta_2,\nu,\sigma_2).
        \notag
    \end{align}
    Assume, for every $b,X,\eta_1,\eta_2,\sigma_1,\sigma_2$, that
    \begin{align}
        C_{A_1,A_2}(b,X;\eta_1,\eta_2,\sigma_1,\sigma_2)
        &=
        C_{B_1,B_2}(b,X;\eta_1,\eta_2,\sigma_1,\sigma_2).
        \notag
    \end{align}
    The following diagram records the three-shared-bond proof subcase used in
    the source, with the external boundary indices restored:
    % Formula: the two inserted two-tensor coefficients agree for every shared
    % bond b and matrix X.  Ink-to-index: in the three-shared-bond proof subcase,
    % X splits the upper contraction, the two north legs are sigma_1 and sigma_2,
    % and the outer stubs expose eta_1 and eta_2.  Boundary signature: both panels
    % have (V,P)=(2,2).  Source verdict: paper_normal.tex:1093--1094 explicitly
    % reduces the arbitrary shared-bond family of eq. (9), lines 1067--1091, to
    % three shared bonds without loss of generality; the theorem's aggregate
    % external virtual boundaries are restored here.
    \begin{center}
        \tenkzkernel
        \begin{tenkz}[rows={wire}, cols=4, bonds=none]
            \tn[at=(1,1), name=twoInjAOne, label pos=s,
                ports={180:virtual:$\eta_1$, 20:virtual, 0:virtual,
                       -20:virtual, 90:physical}]{A_1}
            \tn[at=(1,4), name=twoInjATwo, label pos=s,
                ports={0:virtual:$\eta_2$, 160:virtual, 180:virtual,
                       200:virtual, 90:physical}]{A_2}
            \tn[at=0.45 n of midway twoInjAOne and twoInjATwo,
                name=twoInjAX, skin=ring, species=operator,
                ports={200:virtual, -20:virtual}]{X}
            \tn[at=0.45 s of midway twoInjAOne and twoInjATwo,
                name=twoInjALow, skin=none]{}
            \tnwire{twoInjAOne.20}{twoInjAX.200}
            \tnwire{twoInjAX.-20}{twoInjATwo.160}
            \tnwire{twoInjAOne.0}{twoInjATwo.180}
            \tnwire{twoInjAOne.-20}{twoInjALow}
            \tnwire{twoInjALow}{twoInjATwo.200}
        \end{tenkz}
        $ = $
        \begin{tenkz}[rows={wire}, cols=4, bonds=none]
            \tn[at=(1,1), name=twoInjBOne, label pos=s,
                ports={180:virtual:$\eta_1$, 20:virtual, 0:virtual,
                       -20:virtual, 90:physical}]{B_1}
            \tn[at=(1,4), name=twoInjBTwo, label pos=s,
                ports={0:virtual:$\eta_2$, 160:virtual, 180:virtual,
                       200:virtual, 90:physical}]{B_2}
            \tn[at=0.45 n of midway twoInjBOne and twoInjBTwo,
                name=twoInjBX, skin=ring, species=operator,
                ports={200:virtual, -20:virtual}]{X}
            \tn[at=0.45 s of midway twoInjBOne and twoInjBTwo,
                name=twoInjBLow, skin=none]{}
            \tnwire{twoInjBOne.20}{twoInjBX.200}
            \tnwire{twoInjBX.-20}{twoInjBTwo.160}
            \tnwire{twoInjBOne.0}{twoInjBTwo.180}
            \tnwire{twoInjBOne.-20}{twoInjBLow}
            \tnwire{twoInjBLow}{twoInjBTwo.200}
        \end{tenkz}
    \end{center}
    Then there exists $\lambda\in\C^\times$ such that
    $A_1=\lambda B_1$ and $A_2=\lambda^{-1}B_2$.
    This is the blueprint form of the two-injective tensor comparison in
    \cite[Section~3]{Molnar2018NormalPEPS}.
\end{theorem}
\begin{proof}
    \leanok
    Putting $X=\Id_{D_b}$ on every shared bond contracts all bonds diagonally and,
    by the insertion equality, equates the two fully contracted states. Writing
    each contracted state as a bipartite operator with the linearly independent
    Schmidt families supplied by injectivity, the operator-Schmidt uniqueness of
    \cite[Section~3]{Molnar2018NormalPEPS} produces an invertible gauge $g$ on the
    shared-bond configurations with
    \begin{align}
        A_1(\eta_1,\mu,\sigma_1)
        &=
        \sum_\nu g_{\mu,\nu}B_1(\eta_1,\nu,\sigma_1),
        \notag\\
        B_2(\eta_2,\nu,\sigma_2)
        &=
        \sum_\mu g_{\mu,\nu}A_2(\eta_2,\mu,\sigma_2).
        \notag
    \end{align}
    Substitute both gauge equations into the insertion equality for the
    matrix unit $X=E_{p,q}$ on a single bond $b$, which frees that bond and
    contracts the others. Separating the linearly independent families $B_1$ and
    $A_2$ gives, for every bond $b$, every pair of endpoints $p,q$, and all
    configurations $\nu,\rho$,
    \begin{align}
        [\rho_b=q]\,g_{\rho^{b\to p},\nu}
        &=
        [\nu_b=p]\,g_{\rho,\nu^{b\to q}}.
        \label{eq:peps_residual_equality}
    \end{align}
    Here $\rho^{b\to p}$ replaces the value of $\rho$ at $b$ by $p$. This is the
    bond-wise form of the residual identity of
    \cite[Section~3]{Molnar2018NormalPEPS}, where it is packaged as the equality
    $\Id_\alpha\otimes Z
        =U\otimes\Id_\gamma
        =W_{\alpha\gamma}\otimes\Id_\beta$ of the
    three residual gauges $Z$, $U$, $W$
    (Theorem~\ref{thm:peps_twoInjectiveGaugeScalarReduction}). Reading
    (\ref{eq:peps_residual_equality}) at the two endpoint choices shows that $g$
    vanishes whenever two configurations differ at a bond and that its diagonal
    value is invariant under changing any single bond. Hence
    $g=\lambda\Id$ for a scalar $\lambda$, and invertibility forces
    $\lambda\neq0$.

    Reinserting $g=\lambda\Id$ into the two gauge equations gives
    $A_1=\lambda B_1$ and $B_2=\lambda A_2$, that is
    $A_2=\lambda^{-1}B_2$.
\end{proof}
